Incumbent pairing---the placement of two sitting members of Congress in the same
new district after redistricting---is one of the most politically contested consequences
of any redistricting cycle. Enacted maps routinely minimise pairing by using incumbent
home address data to draw protective boundaries; algorithmic maps, by definition,
cannot. This paper derives the analytical expectation for incumbent pairing under
random redistricting for North Carolina ($k=14$, $n=13$), Wisconsin ($k=8$, $n=8$),
and Texas ($k=38$, $n=36$) under 2020 Census data, and tests whether \textsc{bisect}
algorithmic maps and enacted 2022 congressional maps differ significantly from that
baseline. We find that \textsc{bisect} produces pairing rates of $0.143^\dagger$ (NC),
$0.125^\dagger$ (WI), and $0.105^\dagger$ (TX), all within one standard deviation of
the geographically adjusted random expectation. Enacted maps produce pairing rates of
0.000 (NC), 0.000 (WI), and 0.026 (TX)---significantly below the random baseline in
all three states ($p < 0.05$ under a one-sided binomial test). We conclude that enacted
maps provide deliberate incumbency protection that requires access to incumbent location
data, while \textsc{bisect} maps produce pairing outcomes that are purely a function of
geography. For redistricting litigation, this distinction establishes that the enacted
maps' lower pairing rates are a redistricting choice, not a geographic necessity.
